\begin{mistake}{2026-04-16}{02}

\begin{problemcontent}
设 \(A = \begin{bmatrix} 1 & -2 & 0 \\ 2 & 1 & 5 \\ 0 & 1 & 1 \end{bmatrix}\)，\(B\) 是三阶矩阵，则满足 \(AB = O\) 的所有的 \(B = \underline{\hspace{3cm}}\)。
\end{problemcontent}

\begin{solutioncontent}
设 \(B = (\beta_1, \beta_2, \beta_3)\)，由 \(AB = O\) 得：
\[
A(\beta_1, \beta_2, \beta_3) = (A\beta_1, A\beta_2, A\beta_3) = (0, 0, 0)
\]
即 \(B\) 的列向量 \(\beta_1, \beta_2, \beta_3\) 均为齐次线性方程组 \(Ax = 0\) 的解。

对系数矩阵 \(A\) 进行初等行变换求基础解系：
\[
A = \begin{bmatrix} 1 & -2 & 0 \\ 2 & 1 & 5 \\ 0 & 1 & 1 \end{bmatrix}
\xrightarrow{r_2 - 2r_1} \begin{bmatrix} 1 & -2 & 0 \\ 0 & 5 & 5 \\ 0 & 1 & 1 \end{bmatrix}
\xrightarrow[r_3 - r_2]{r_2 \div 5} \begin{bmatrix} 1 & -2 & 0 \\ 0 & 1 & 1 \\ 0 & 0 & 0 \end{bmatrix}
\xrightarrow{r_1 + 2r_2} \begin{bmatrix} 1 & 0 & 2 \\ 0 & 1 & 1 \\ 0 & 0 & 0 \end{bmatrix}
\]
对应的同解方程组为 \(x_1 = -2x_3,\ x_2 = -x_3\)。
取自由变量 \(x_3 = 1\)，得到 \(Ax = 0\) 的基础解系：
\[
\xi = \begin{bmatrix} -2 \\ -1 \\ 1 \end{bmatrix}
\]
故方程组的解均可表示为 \(k\xi\)。因为 \(\beta_i\) 都是方程组的解，设 \(\beta_i = k_i\xi \ (i=1,2,3)\)，则有：
\[
B = (k_1\xi, k_2\xi, k_3\xi) = \xi \begin{bmatrix} k_1 & k_2 & k_3 \end{bmatrix} = \begin{bmatrix} -2 \\ -1 \\ 1 \end{bmatrix} \begin{bmatrix} k_1 & k_2 & k_3 \end{bmatrix} = \begin{bmatrix} -2k_1 & -2k_2 & -2k_3 \\ -k_1 & -k_2 & -k_3 \\ k_1 & k_2 & k_3 \end{bmatrix}
\]
其中 \(k_1, k_2, k_3\) 为任意常数。
\end{solutioncontent}

\mistakeknowledge{
\begin{itemize}
 \item \textbf{核心考点}：矩阵分块乘法 \(AB = A(\beta_1, \cdots, \beta_n) = (A\beta_1, \cdots, A\beta_n)\)，将矩阵方程转化为求解线性方程组。
 \item \textbf{易错点}：求解通解后，误将 \(B\) 的三列写成完全相同的参数（如每一列都是 \(k\)），忽略了矩阵 \(B\) 的每一列解是相互独立的，必须用不同的任意常数 \(k_1, k_2, k_3\) 表示。
 \item \textbf{关联知识}：初等行变换求行最简形矩阵；齐次线性方程组基础解系的构造。
\end{itemize}}

\end{mistake}
